% ============================
% LaTeX Template - Problem Set
% ============================
\documentclass[11pt]{article}

% ---------- Page + layout ----------
\usepackage[margin=1in]{geometry}
\usepackage{setspace}
\setstretch{1.1}
\usepackage{parskip}

% ---------- Math ----------
\usepackage{amsmath, amssymb, amsfonts, amsthm, mathtools}
\usepackage{threeparttable}
\usepackage[utf8]{inputenc}

% ---------- Tables ----------
\usepackage{booktabs}

% ---------- Links ----------
\usepackage[hidelinks]{hyperref}

% ---------- Header / footer ----------
\usepackage{fancyhdr}
\pagestyle{fancy}
\fancyhf{}
\lhead{\textbf{\course}}
\chead{\textbf{\pset}}
\rhead{\textbf{\name}}
\cfoot{\thepage}

% =========================
% Metadata (edit if needed)
% =========================
\newcommand{\name}{Tate Mason}
\newcommand{\course}{Econometrics II}
\newcommand{\pset}{Problem Set 7}
\newcommand{\duedate}{April 8, 2026}

% =========================
% Question / solution macros
% =========================
\newcounter{q}
\renewcommand{\theq}{\arabic{q}}

\newcommand{\question}[1]{%
  \refstepcounter{q}
  \vspace{0.75em}
  \noindent{\Large\textbf{Question \theq.} \normalsize\textbf{#1}}%
  \vspace{0.25em}
  \par
}

\newcounter{subq}[q]
\renewcommand{\thesubq}{\alph{subq}}

\newcommand{\subquestion}[1]{%
  \refstepcounter{subq}
  \vspace{0.35em}
  \noindent\textbf{(\thesubq)} \textbf{#1}\par
}

\newenvironment{solution}{%
  \vspace{0.15em}
  \noindent\textit{Solution. }\ignorespaces
}{\vspace{0.5em}\par}

% =========================
% Easy math macros
% =========================
\newcommand{\R}{\mathbb{R}}
\newcommand{\N}{\mathbb{N}}
\DeclareMathOperator{\E}{\mathbb{E}}
\DeclareMathOperator{\Var}{Var}
\DeclareMathOperator{\Cov}{Cov}
\DeclareMathOperator{\plim}{plim}
\DeclareMathOperator*{\argmin}{argmin}
\DeclareMathOperator*{\argmax}{argmax}
\newcommand{\P}{\mathbb{P}}

\newcommand{\vct}[1]{\mathbf{#1}}
\newcommand{\mtx}[1]{\mathbf{#1}}
\newcommand{\dd}{\,\mathrm{d}}
\newcommand{\eps}{\varepsilon}
\newcommand{\1}{\mathbf{1}}

% =========================
% Document
% =========================
\begin{document}

% ---------- Title block ----------
{\Large\textbf{\course} \hfill \textbf{\pset}}\\
\textbf{Name:} \name \hfill \textbf{Due:} \duedate\\
\rule{\textwidth}{0.4pt}

\question{Use \texttt{dataHW7\_Problem1.mat} to complete this exercise. The data is identical to
\texttt{dataHW6\_Problem2.mat} --- meaning the structure of the problem is the same as are the
underlying parameters. The purpose of this problem is to estimate a simple dynamic discrete choice
problem using CCPs. The data contain information on the choices ($Y$'s), characteristics that do
not vary over time ($X_1$'s), and characteristics that do vary over time ($X_{1t}$'s). The context
of the problem is a dynamic model of labor supply. $Y$ takes values $\{0, 1, 2, 3\}$ reflecting
retirement, flex-time, part-time, or full-time participation. The first column of $X_1$ allows for
choice-specific constants, while the second column is an indicator for gender. $X_{1t}$ is a
time-varying measure of health.
Both sets of $X$'s are individual-level characteristics, so the coefficients on these variables must
vary by choice. There is a cost to switching labor supply that is common across all switches. The
choice in period 0 is given by $LY$. Once an individual retires (chooses option 0) they make no
further choices: if $Y_{it} = 0$ then $Y_{it+1} = 0$. The columns of $Y$ and $X_{1t}$ indicate
observations over time, ten periods in all, so that period 10 is a static problem. Individuals are
uncertain about future health. Given $X_{1t}$ and $Y_t$, the health transition is:
%
\begin{equation}
    X_{1,t+1} = \alpha_0 + \alpha_1 X_{1t} + \alpha_2 \mathbf{1}(Y_t = 2)
                + \alpha_3 \mathbf{1}(Y_t = 3) + \varepsilon_{t+1}
\end{equation}
%
for $Y_t > 0$, where $\varepsilon_{t+1} \sim \mathcal{N}(0, \sigma)$ is unknown to the individual
until time $t+1$.}

\subquestion{Estimate the transition parameterts, the $\alpha$'s and $\sigma$.}

\begin{solution}
  \begin{table}[htbp]
  \centering
    \caption{OLS Estimates}
    \begin{tabular}{lc}
      \hline
      & Parameter \\
      \hline
      $\alpha_{0}$    & 0.015 \\
      $\alpha_{1}$ & -0.754 \\
      $\alpha_{pt}$ & 0.035 \\
      $\alpha_{ft}$ & -0.257 \\
      \hline
      $\sigma$ & 1.003 \\
      \bottomrule
    \end{tabular}
  \end{table}
\end{solution}

\clearpage

\subquestion{Estimate the CCP's via a single logit model of retirement based on LY, X1, X1t, and a set of time dummies.}

\begin{solution}
  \begin{table}[h!]
    \centering
    \begin{threeparttable}
    \caption{Part B: Conditional Choice Probability Logit Estimates}
      \begin{tabular}{lc}
        \toprule
        Variable & Coef. \\
        \midrule
        $X_1$ (Constant)                           & $-$3.6991  \\
        $X_2$ (Gender)                             & $-$0.8556  \\
        $X_t$ (Health)                             & $-$0.3494  \\
        $Y_0$ (Initial state)                      &    0.0115  \\
        $\mathbf{1}(\bar{Y}_{t-1}=2)$ (Part-time) & $-$0.1116\\
        $\mathbf{1}(\bar{Y}_{t-1}=3)$ (Full-time) & $-$0.0807 \\
        $t=2$  & $-$0.0443  \\
        $t=3$  &    0.1227  \\
        $t=4$  &    0.2846  \\
        $t=5$  &    0.6496  \\
        $t=6$  &    1.0620  \\
        $t=7$  &    1.5239  \\
        $t=8$  &    2.1992  \\
        $t=9$  &    3.6510  \\
        \midrule
        Observations & \multicolumn{4}{l}{39,413} \\
        Pseudo $R^2$ & \multicolumn{4}{l}{0.2354} \\
        \bottomrule
      \end{tabular}
    \end{threeparttable}
  \end{table}
\end{solution}

\subquestion{Use the CCP estimates and the future value terms calculated in (C) to estimate choice specific parameters, switching costs, and the discount factor via multinomial logit.}

\begin{solution}
\begin{table}[h!]
  \centering
  \begin{threeparttable}
  \caption{Part D: Dynamic Discrete Choice Estimates (CCP Method)}
    \begin{tabular}{lcc}
      \toprule
      Parameter & Estimate & Truth \\
      \midrule
      \multicolumn{4}{l}{\textit{Flex-time ($j=1$)}} \\
      \quad Constant & $-$0.8940 & $-$0.75 \\
      \quad Gender   &    0.5936 &    0.25 \\
      \quad Health   &    0.1695 &    0.25 \\
      \midrule
      \multicolumn{4}{l}{\textit{Part-time ($j=2$)}} \\
      \quad Constant & $-$0.3305 & $-$0.25 \\
      \quad Gender   &    0.8890 &    0.50 \\
      \quad Health   &    0.4162 &    0.50 \\
      \midrule
      \multicolumn{4}{l}{\textit{Full-time ($j=3$)}} \\
      \quad Constant & $-$1.0477  & $-$0.85 \\
      \quad Gender   &    1.1255  &    0.75 \\
      \quad Health   &    0.6673  &    0.75 \\
      \midrule
      Switching cost          & 0.4239 & 0.40 \\
      Discount factor $\beta$ & 0.7796 & 0.80 \\
      \bottomrule
    \end{tabular}
  \end{threeparttable}
\end{table}
\end{solution}

\end{document}
